\normalsize{El apartado ``Requisitos de las acreditaciones internacionales EUR-ACE y ABET'' de la Normativa de TFT de la ETSIT-UPM establece que ``\textit{La memoria del TFT del GITST, GIB y MUIT, y en general la de aquellas titulaciones que hayan obtenido o para las que se desee solicitar una acreditación internacional EUR-ACE o ABET, debe mostrar conciencia de la responsabilidad de la aplicación práctica de la ingeniería, el impacto social y ambiental, el compromiso con la ética profesional, la responsabilidad y las normas de la aplicación práctica de la ingeniería, así como sobre las prácticas de gestión de proyectos, gestión, y control de riesgos, entendiendo sus limitaciones}''.}
\vspace{0.1cm}

\section*{A1. Introducci\'on}
\normalsize{This project develops and evaluates a machine-learning system
that screens for Parkinson's disease from contactless radar recordings of
gait. It touches three sensitive domains at once: \textbf{health data} from
a vulnerable population, \textbf{clinical claims} produced by
machine-learning models, and a sensing technology whose reach into private
spaces is precisely its selling point. The relevant dimensions are
therefore ethical (validity of clinical claims, privacy), social (access to
earlier screening for an ageing population), economic (low-cost monitoring
against expensive laboratory instrumentation) and, to a lesser degree,
environmental (computational footprint).}
\vspace{0.3cm}

\section*{A2. Descripci\'on de impactos relevantes relacionados con el proyecto}
\normalsize{Three impacts dominate, each with its stakeholders:

\begin{itemize}
  \item \textbf{The validity of clinical machine-learning claims.}
        Stakeholders: patients, clinicians, the research field. An
        overstated screening accuracy harms twice: false alarms send
        healthy people to neurology consultations, and false trust delays
        real diagnoses. This thesis found that the field's evaluation
        habits inflate performance substantially, and made honest
        evaluation its central contribution.
  \item \textbf{Privacy of gait data.} Stakeholders: monitored persons and
        their households. Radar records no image, which is its privacy
        advantage; yet this work also measured that a person's walk
        identifies them (all 58 of 58 subjects distinguishable), so gait
        recordings are biometric data under the GDPR and must be treated as
        such: anonymised identifiers, restricted access, purpose
        limitation. The dataset used here is fully anonymised and was
        collected under the clinical validation study of the source
        system.
  \item \textbf{Access and cost of motor assessment.} Stakeholders: health
        systems, patients in areas without movement-disorder specialists. A
        radar Timed-Up-and-Go, the one measurement this thesis validates as
        trustworthy, is cheap, contactless and automatable, and could
        extend objective motor assessment beyond specialist clinics.
\end{itemize}}
\vspace{0.1cm}

\section*{A3. An\'alisis detallado de alguno de los impactos}
\normalsize{The impact analysed in depth is the first: \textbf{what happens
when clinical machine-learning performance is overstated}, because this
thesis measured the mechanism directly. On this very dataset, the
evaluation design used by most published work produces subject-level
scores that look clinic-ready; scoring the same models on people excluded
from training returns them to a band indistinguishable from a predictor
that knows only how long the recording lasted. A deployment decision taken
on the first number would field a screening tool whose real behaviour on
new patients is close to a biased guess: when the field's own network was
scored on unseen people (Section~\ref{sec:disc-alexnet}), catching
21 of 25 patients cost falsely flagging 18 of 33 healthy controls. The
ethical posture adopted is therefore procedural rather than declarative:
pre-registration of the analysis plan, subject-independent validation,
explicit null floors beside every reported score, confidence intervals at
the subject level, and publication of the negative result. This is the
engineering-responsibility standard the EUR-ACE criteria ask the graduate
to demonstrate, applied to the claim itself rather than to the artefact.}
\vspace{0.3cm}

\section*{A4. Conclusiones}
\normalsize{The project's sustainability profile is favourable: it consumes
only modest computation (a single personal computer; no data-centre
training), reuses an existing clinical dataset rather than collecting a new
one, and its product is methodological, a protocol that prevents overstated
clinical claims. The sustainability criteria added genuine value: treating
the validity of the claim as the primary ethical object is what turned an
apparently negative result into the thesis's contribution, and the privacy
analysis (gait as a biometric) is itself one of the work's findings.}
